\SourceLeaf{27}
\textsc{Gauss}' eigene Ausführungen beziehen sich auf eine spezielle Einteilung der
$(p-1)$ Reste mod.\ $p$ in drei Klassen $S$, $S'$, $S''$, wobei die
Darstellung von $p$ in der Form $(mm+3nn)$ verwendet wird. Vermutlich haben
aber die Schwierigkeiten, die mit dieser Bestimmung der Zahl $f$ verbunden
sind, \textsc{Gauss} davon abgehalten, den hier eingeschlagenen Weg weiterzuverfolgen.
Gerade an dieser Stelle wird die Notwendigkeit besonders deutlich, den
Zahlbereich auf die komplexen ganzen Zahlen $(a+b\rho)$ zu erweitern, wobei an
die Stelle von $f$ die Einheitswurzel
$\rho = \frac{-1+i\sqrt{3}}{2}$ selbst tritt.

Die fraglichen Untersuchungen wurden jedenfalls vor 1809 durchgeführt und
schlossen wahrscheinlich unmittelbar an die genannten Entwicklungen über
quadratische Reste an, die im Januar 1808 der Gött. Gesellsch. d. Wiss. vorgelegt
wurden. Nach Notizen in einem Tagebuch, das \textsc{Gauss} seit Ende März 1796
über viele Jahre hinweg über seine Untersuchungen führte\HistoricalNote{\textsuperscript{*)}}{*)}{Über dieses höchst wertvolle Tagebuch sollen in
Bd.\ 9 der ges. Werke nähere Angaben gemacht werden.}, hat er sich Mitte Februar
1807 eingehender mit der Theorie der kubischen und biquadratischen Reste befasst.
Wahrscheinlich schloss sich daran auch unmittelbar der Ersatz von $f$ durch
$\rho$ an, wie aus der Urschrift der oben mitgeteilten Notizen hervorgeht. Man
darf demnach annehmen, dass \textsc{Gauss} bereits um die angegebene Zeit (1808) den
so folgenreichen Schritt der Einführung der komplexen ganzen Zahlen endgültig
vollzogen hat.

\textsc{Gauss} hat auch für die komplexen Primzahlen $2-\rho$ und $3-\rho$ die drei
oben wiedergegebenen Klassen $S$, $S'$, $S''$ aufgestellt, ebenso für den
Primteiler $3-2\rho$ von 19. Besonders interessant ist, dass \textsc{Gauss} diese Überlegungen,
wie entsprechend auch bei den biquadratischen Resten, in geometrische Form
brachte. Die komplexen ganzen Zahlen $a+b\rho$ liefern Punkte der Ebene, die ein
rhombisches Punktgitter bilden. Die Gitterpunkte eines einzelnen Systems
$S$ lassen sich dabei auf einen Bereich beschränken, der im neueren Sinne
als Diskontinuitätsbereich einer unendlichen Gruppe linearer Substitutionen
aufgefasst werden kann. \textsc{Gauss} hat für verschiedene Fälle Zeichnungen
angefertigt. Die für die Primzahl $6-\rho$ ist oben wiedergegeben. Es würde sich
hier um die Gruppe
\[
u' = \rho^{\nu} u + (a+b\rho),
\]
handeln, wobei $u$ eine komplexe Variable bezeichnet, $\nu = 1, 2, 3$ zu setzen ist
und $a+b\rho$ alle komplexen ganzen Zahlen durchläuft, die der Kongruenz
\[
a+b\rho \equiv 0 \quad(\mathrm{mod.}\ 6-\rho)
\]
genügen. Die in der \textsc{Gauss}'schen Figur angegebenen Bereiche beziehen sich
auf die Gruppe, die durch Hinzunahme der Substitution $u' = \pm u$ erweitert wird.
Die freie Wahl des Repräsentantensystems $S$ entspricht nun der freien Wahl
der Bereichsgrenzen. Welche Absicht \textsc{Gauss} mit der von ihm gewählten Gestalt
dieser Grenzen verfolgte, ist leider nicht angegeben.

Die zwölf Thesen, die am Schluss der obigen Notizen unter [2] zusammengestellt
sind, folgen in der Urschrift direkt auf die Betrachtungen der Zahlen $a+b\rho$.
Das System aller komplexen ganzen Zahlen $a+b\rho$ erscheint hier eindeutig ersetzt
durch das mod.\ $(1+x+xx)$ reduzierte System aller Polynome in $x$ mit rationalen ganzen Zahlen
als Koeffizienten.
\par\nobreak\noindent\hfill\textsc{Fricke}.\par
\endgroup
